\section{Objetivos}

\subsection{Objetivo general}

Optimizar y evaluar el clasificador de Doctor Maíz sobre el corpus actualizado,
con un protocolo reproducible que permita seleccionar una alternativa
experimental, medir su comportamiento por clase y condición de captura, y
comprobar su integración en un prototipo de diagnóstico asistido.

\subsection{Objetivos específicos}

\begin{enumerate}
  \item Reconstruir el dataset definitivo de la etapa, validar sus archivos,
  eliminar duplicados exactos y congelar una división que mantenga separado el
  holdout final.

  \item Optimizar una cabeza clasificadora sobre EfficientNet-B0 mediante
  búsqueda TPE y pruning con Optuna, usando macro-F1 de validación como función
  objetivo y sin consultar el holdout.

  \item Comparar EfficientNet-B0, ShuffleNetV2-x1.0, MobileNetV3 Small y
  FastViT-T8 bajo el mismo preprocesamiento, partición y configuración de
  cabeza, incluyendo calidad, cantidad de parámetros y latencia CPU del
  backbone.

  \item Medir complementariedad de errores y contrastar el mejor modelo
  individual con soft voting uniforme y ponderado. Los pesos deben aprenderse
  únicamente en validación.

  \item Estimar la estabilidad de la selección con cinco folds estratificados
  por clase y entorno, y efectuar una sola evaluación sobre el holdout
  congelado con métricas globales y por clase.

  \item Revisar el principal cuello de botella de la Etapa 1,
  \texttt{potassium\_deficiency}, cuantificando sus confusiones con nitrógeno,
  fósforo y las demás categorías.

  \item Explicar casos finales con Grad-CAM y LIME, y describir brechas de
  precision, recall y F1 asociadas con clase, laboratorio/campo, fuente,
  resolución y nitidez.

  \item Verificar el flujo real de la aplicación móvil y producir una demo
  ejecutable que muestre entrada, preprocesamiento, Top-3, confianza, latencia
  y advertencia de uso.

  \item Traducir las limitaciones encontradas en condiciones de uso,
  recomendaciones para extensionistas y propuestas concretas de política
  pública, sin presentar beneficios no medidos.
\end{enumerate}

Estos objetivos combinan tres niveles que en la primera etapa todavía estaban
separados: calidad estadística, comportamiento bajo condiciones distintas y
operación dentro de una herramienta usable. Una mejora pequeña en macro-F1 no
compensa una caída fuerte en campo; del mismo modo, una aplicación bien
diseñada no corrige una clase minoritaria mal reconocida. La selección final se
lee con las tres dimensiones.
